\chapter{Sample Chapter}\label{ch:1-sample}

% Opening motivation: a concrete engineering problem the chapter will solve.

\begin{objectives}
After studying this chapter, you should be able to:
\begin{enumerate}[label=\arabic*.]
  \item compute \dots
  \item explain why \dots
\end{enumerate}
\end{objectives}

\section{First Concept}\label{sec:1-first}
Why, intuition, model, formal statement, check. Cite the literature the chapter relies on, e.g.\ \citet{boyce2014}.
\begin{equation}\label{eq:1-ohm}
  V = R\,I
\end{equation}

\begin{example}[Current through a heater]\label{ex:1-heater}
\exgiven A heating element of resistance \SI{24}{\ohm} is connected to \SI{230}{\volt}.
\exfind The current drawn.
\exapproach Apply \eqref{eq:1-ohm} directly; assume a purely resistive load.
\exsolution $I = V/R = \SI{230}{\volt}/\SI{24}{\ohm} \approx \SI{9.6}{\ampere}$.
\excheck Units: \si{\volt\per\ohm} = \si{\ampere}; magnitude typical of a \SI{2.2}{\kilo\watt} heater.
\exdiscussion A \SI{10}{\ampere} circuit would be marginal; a \SI{16}{\ampere} circuit is appropriate.
\end{example}

\begin{practicenote}
Real heater resistance rises with temperature, so the cold inrush current exceeds the running current.
\end{practicenote}

\section*{\summaryname}
Key ideas and equations (see Example~\ref{ex:1-heater}).

\section*{Problems}
\begin{problems}
  \item A basic problem.
  \item\probint{} An intermediate problem.
  \item\probdesign{} A design problem.
\end{problems}

% ---- End-of-chapter references (books and peer-reviewed papers only) ----
\bibliographystyle{plainnat}
\bibliography{references}
